\documentclass[10pt,a4paper]{amsart}
\usepackage[margin=19mm]{geometry}
\usepackage{amsmath,amssymb,mathtools}
\usepackage[scr=rsfs,cal=euler]{mathalfa}
\usepackage{xfrac}
\usepackage[unicode,hidelinks,bookmarks=false]{hyperref}
\newcommand{\ie}{i.e.\ }
\newcommand{\cf}{cf.\ }
\newcommand{\resp}{respectively\ }
\newcommand{\Subsection}{\subsection}
\providecommand{\nobreakdash}{\nobreak\hbox{-}}
\setlength{\emergencystretch}{2em}
\multlinegap0pt
\usepackage{float}
\usepackage{amssymb}
\usepackage[shortlabels]{enumitem}
\newtheorem{theorem}{Theorem}
\newtheorem{definition}{Definition}
\newtheorem{lemma}{Lemma}
\newcommand\Gbb{{\mathbb{G}}}
\newcommand\Hol{{\mathrm{Hol}}}
\newcommand\vol{\mathrm{vol}}
\title{Universality of two-dimensional Markovian holonomy fields}
\author{Thibaut Lemoine, Elias Nohra}
\date{Source: arXiv:2608.15923v1 (CC BY 4.0)}
\begin{document}
\maketitle

\noindent\emph{Source and licence.} Universality of two-dimensional Markovian holonomy fields. Version arXiv:2608.15923v1 (CC BY 4.0), distributed by arXiv under the Creative Commons Attribution 4.0 International licence, http://creativecommons.org/licenses/by/4.0/, as recorded in the arXiv licence metadata for this exact version and shown on the abstract page https://arxiv.org/abs/2608.15923v1. Full licence notice: \texttt{licenses/arxiv-2608.15923-CC-BY-4.0.txt}.\par\smallskip

\noindent\textit{Editorial scope.} Counted unit: the article's theorem thm:Fourier\_PF (source0.tex lines 251--265). The article's complete original proof of it is delivered with it (source0.tex lines 711--812, one \texttt{proof} environment of 102 lines, which the article places away from the statement). No mathematics is written, repaired or re-derived: the statement and the proof are the article's own lines, in the article's own notation, and no retained line is substituted or re-flowed.  Retained uncounted as the source-local context the proof needs: lines 446--484; lines 647--651; lines 653--665; lines 682--697.  One declared adaptation: exactly 1 ancillary strings inside the retained spans are omitted (image-inclusion commands and, where the source repeats a label, the repeated label definitions) and no other line differs from the source; the figure environments, with their placement, captions and labels, are kept, and the package records the omitted strings in fidelity receipts bound to the affected bodies.  No retained line contains a citation, so the wrapper carries no bibliography and no bibliography entry.  The retained lines contain the article's own diagram code, which the wrapper typesets with the standard tikz packages; no image file is delivered and the package declares no asset dependency.  Editorial formatting only, in addition: an AMS \texttt{amsart} preamble that restates the article's own declarations and macros used by the retained lines, namely the environments \texttt{theorem}, \texttt{definition}, \texttt{lemma} on separate counters, together with the standard packages those lines need.  The retained environments print in this document as Theorem 1, Definition 1, Lemma 1 in order of appearance, whereas the article prints the same items with the numbering of its own full document.  The complete native source of the article as distributed in the arXiv e-print for this exact version is bundled unchanged as \texttt{sources/207795/source0.tex}.
\begin{definition}\label{def:discrete-YM-measure}
Let $\Sigma$ be a compact surface of genus $g$ with boundary components
$b_1,\ldots,b_n$, endowed with an area measure $\vol$. Let
$\Gbb=(V,E,F)$ be a topological map embedded in $\Sigma$, let $G$ be a
compact group, and let $(Q_t)_{t>0}$ be a family of central continuous
probability densities on $G$. Fix conjugacy classes
$[t_1],\ldots,[t_n]$ and write
$\mathfrak t=([t_1],\ldots,[t_n])$.

For every $i$, choose an oriented edge $e_i$ contained in the $i$-th
boundary component; these edges are pairwise distinct. Once all variables
except $U_{e_i}$ are fixed, the boundary holonomy has the form
\[
 \Hol_{\partial b_i}(U)=a_i(U)U_{e_i}^{\varepsilon_i}b_i(U),
 \qquad \varepsilon_i\in\{\pm1\},
\]
so, for every $y_i\in G$, there is a unique value of $U_{e_i}$ for which
\[
 \Hol_{\partial b_i}(U)=y_it_iy_i^{-1}.
\]
Denote by $U^{\mathbf y}$ the resulting connection, where
$\mathbf y=(y_1,\ldots,y_n)$. The \emph{discrete Yang--Mills measure}
with weight $Q$ and boundary condition $\mathfrak t$ is the finite measure
characterized by
\begin{equation}\label{eq:DS-rigorous}
\begin{aligned}
 \int_{\Omega^1(\Gbb,G)}\! \Phi(U)\,
 d\mu_{\Gbb,G,Q,\vol,\mathfrak t}(U)
 :=
 \int_{G^{E\setminus\{e_1,\ldots,e_n\}}}\int_{G^n}
 &\Phi(U^{\mathbf y})
 \prod_{f\in F}Q_{|f|}\bigl(\Hol_{\partial f}(U^{\mathbf y})\bigr)
 \,d\mathbf y\,dU
\end{aligned}
\end{equation}
for every bounded measurable $\Phi$. This definition is independent of the
chosen boundary edges, by invariance of Haar measure under left and right
translations and inversion.
\end{definition}

\begin{equation}\label{eq:character-face-gluing}
 \int_G\chi_\lambda(ax)\chi_\mu(x^{-1}b)\,dx
 =
 \delta_{\lambda\mu}\frac{\chi_\lambda(ab)}{d_\lambda}.
\end{equation}

\begin{lemma}[Tree gauge fixing]\label{lem:tree-gauge-fixing}
Let $\Gamma=(V,E)$ be a finite connected graph, let $T\subset E$ be a
spanning tree, and fix a root $o\in V$. If $\Phi:G^E\to\mathbb C$ is
integrable and invariant under the gauge action of $G^V$, then
\[
 \int_{G^E}\Phi(U)\,dU
 =
 \int_{G^{E\setminus T}}\Phi(U^T)\,dU,
\]
where $U^T_e=1_G$ for $e\in T$ and $U^T_e=U_e$ otherwise.
The same statement holds for the boundary-conditioned integrals of
Definition~\ref{def:discrete-YM-measure}.
\end{lemma}

\begin{lemma}[One-face normal form]\label{lem:one-face-normal-form}
Let $\mathbb G_0$ be a connected one-face topological map on the compact
oriented surface $\Sigma_{g,n}$ and suppose that every boundary component is
carried by $\mathbb G_0$. After contracting a spanning tree, the resulting
one-vertex ribbon graph has $2g+n$ oriented loop edges. There is a free basis
\[
 a_1,b_1,\ldots,a_g,b_g,c_1,\ldots,c_n
\]
of its edge-path group such that, with the face and boundary orientations induced by the orientation of the surface, the unique face word is, up to cyclic
permutation,
\[
 \prod_{i=1}^g[a_i,b_i]\prod_{j=1}^n c_j,
\]
and $c_j$ represents the $j$-th oriented boundary component. The
corresponding change of edge variables preserves product Haar measure.
\end{lemma}

\begin{theorem}[Partition function]\label{thm:Fourier_PF}
Let $\Sigma_{g,n}$ be a compact connected oriented surface and let $\Gbb=(V,E,F)$ be a topological map embedded in $\Sigma_{g,n}$, with every boundary component carried by $\Gbb$. Let $(Q_t)$ be a family of central continuous probability densities and assume the graph-level spectral summability condition
\begin{equation}\label{eq:PF-spectral-summability}
 \sum_{\lambda\in\widehat G}
 d_\lambda^{2-2g}
 \prod_{f\in F}|a_{|f|}(\lambda)|<\infty.
\end{equation}
Then the partition function has the absolutely convergent expansion
\begin{equation}\label{eq:pf-state-sum-intro}
 Z_{\Gbb,G,Q,\vol}(t_1,\ldots,t_n)
 =\sum_{\lambda\in\widehat G}
 \left(\prod_{f\in F}a_{|f|}(\lambda)\right)
 d_\lambda^{2-2g}\prod_{j=1}^n\frac{\chi_\lambda(t_j)}{d_\lambda}.
\end{equation}
\end{theorem}

\begin{proof}[Proof of Theorem~\ref{thm:Fourier_PF}]
Write $F=\{f_1,\ldots,f_q\}$. For $\varepsilon>0$, regularize each
plaquette density by the heat kernel:
\[
 Q_{|f|}^{(\varepsilon)}
 :=Q_{|f|}*p_\varepsilon.
\]
These are central smooth probability densities, hence Peter--Weyl
regular, and their normalized Fourier coefficients are
\[
 a_{|f|}^{(\varepsilon)}(\lambda)
 =e^{-\frac{\varepsilon}{2}c_2(\lambda)}
 a_{|f|}(\lambda).
\]
Let $Z_\varepsilon$ denote the partition function obtained by replacing
every face weight by $Q_{|f|}^{(\varepsilon)}$. We may multiply the
absolutely and uniformly convergent character expansions and integrate
term by term.

Consider the dual adjacency graph whose vertices are the faces of
$\mathbb G$ and whose edges correspond to internal primal edges separating
two distinct faces. This dual graph is connected: a
path in the interior of the connected surface joining points in two
faces, chosen transverse to the graph and avoiding its vertices, gives a
chain of adjacent faces. Choose a spanning tree $T^*$ of this dual graph.
Integrate successively the $q-1$ primal edge variables dual to $T^*$. An example is displayed in Figure~\ref{fig:tree-integration}. At
each step, Schur orthogonality in the form
\eqref{eq:character-face-gluing} forces the two incident face labels to
agree, merges the character words, and contributes one factor
$d_\lambda^{-1}$. Since $T^*$ connects all faces, all labels are equal to
a single $\lambda$. The $q$ factors $d_\lambda$ from the character
expansions and the $q-1$ gluing factors leave one factor $d_\lambda$. We
obtain
\begin{equation}\label{eq:PF-one-face-intermediate}
 Z_\varepsilon(t_1,\ldots,t_n)
 =
 \sum_{\lambda\in\widehat G}
 e^{-\frac{q\varepsilon}{2}c_2(\lambda)}
 \left(\prod_{f\in F}a_{|f|}(\lambda)\right)d_\lambda I_\lambda,
\end{equation}
where $I_\lambda$ is the boundary-conditioned integral of
$\chi_\lambda$ of the boundary word of the unique remaining face.

\begin{figure}[h!]
 \centering
 
 \caption{An example of the successive integration for a topological map in $\Sigma_{2,0}$. The dual graph is dashed, and the dual tree is in red. The oriented edges are glued pairwise. In the first step, the whole dual graph is displayed, but all non-tree dual edges are omitted in the next steps. Each step consists in integrating the corresponding primal edge (in bold) to an edge of the dual tree, thereby merging two faces of the primal graph.}
 \label{fig:tree-integration}
\end{figure}

The primal graph remaining after these deletions is connected and has one
face. Apply Lemma~\ref{lem:tree-gauge-fixing} to a spanning tree and then
Lemma~\ref{lem:one-face-normal-form}. Thus
\[
 I_\lambda
 =
 \int_{G^{2g+n}}
 \chi_\lambda\!\left(
 \prod_{i=1}^g[A_i,B_i]\prod_{j=1}^n C_j
 \right)
 \prod_{j=1}^n\delta_{[t_j]}(dC_j)
 \prod_{i=1}^g dA_i\,dB_i.
\]
Using orbital coordinates $C_j=U_jt_jU_j^{-1}$ and Schur orthogonality,
\[
 \int_{G^2}\chi_\lambda([A,B]x)\,dA\,dB
 =\frac{\chi_\lambda(x)}{d_\lambda^2},
 \qquad
 \int_G\chi_\lambda(xUtU^{-1})\,dU
 =\frac{\chi_\lambda(x)\chi_\lambda(t)}{d_\lambda}.
\]
Iterating these identities yields
\[
 I_\lambda
 =d_\lambda^{-2g+1-n}
 \prod_{j=1}^n\chi_\lambda(t_j).
\]
Substitution in~\eqref{eq:PF-one-face-intermediate} gives
\begin{align*}
 Z_\varepsilon(t_1,\ldots,t_n)
 &=
 \sum_{\lambda\in\widehat G}
 e^{-\frac{q\varepsilon}{2}c_2(\lambda)}
 \left(\prod_{f\in F}a_{|f|}(\lambda)\right)
 d_\lambda^{2-2g}
 \prod_{j=1}^n\frac{\chi_\lambda(t_j)}{d_\lambda}.
\end{align*}

Since $Q_{|f|}*p_\varepsilon\to Q_{|f|}$ uniformly on $G$ for each
face, the finite product of plaquette densities converges uniformly and
$Z_\varepsilon\to Z_{\mathbb G,G,Q,\vol}$. On the series side,
$|\chi_\lambda(t_j)|\leq d_\lambda$ and
$e^{-q\varepsilon c_2(\lambda)/2}\leq1$, so the absolute value of the
$\lambda$-term is bounded by
\[
 d_\lambda^{2-2g}
 \prod_{f\in F}|a_{|f|}(\lambda)|,
\]
which is summable by~\eqref{eq:PF-spectral-summability}. Dominated
convergence as $\varepsilon\downarrow0$ proves
\eqref{eq:pf-state-sum-intro} and its absolute convergence.
\end{proof}

\end{document}
